% ============================================================
%  chapter/intro.tex — Getting started
%
%  This chapter doubles as a mini-manual.  Write your real first
%  chapter in this file, or add new files in chapter/ and pull them
%  in from main.tex with \include{chapter/yourfile}.
% ============================================================
\chapter{Getting Started}\label{ch:intro}

Welcome to your book. Everything you see on these pages is sample
content that exists to show off the design and to stress-test the
compiler. Replace it freely.

\section{Writing text}
Ordinary paragraphs need no markup at all: just write, and leave one
blank line between paragraphs. \emph{Italics}, \textbf{bold} and
\textsc{small caps} are inline commands; footnotes go at the bottom of
the page.\footnote{Like this one. Footnotes are numbered per chapter.}

Quotation marks are best typed with the \texttt{\textbackslash enquote} command from
\emph{csquotes}, which adapts to the document language:
\enquote{like this}, rather than by typing quote characters by hand.

\subsection{Lists}
Three kinds come built in, tuned by the \emph{enumitem} package:
\begin{itemize}[itemsep=2pt]
  \item an itemised list,
  \item like this one;
  \item nesting works to any sane depth.
\end{itemize}
\begin{enumerate}[itemsep=2pt]
  \item an enumerated list,
  \item numbered automatically;
  \item cross-referenced anywhere as (\ref{enum:start}).
        \label{enum:start}
\end{enumerate}

\section{Cross-references}
Every chapter, section, figure, table and equation can carry a
\texttt{\textbackslash label}; refer to it with \texttt{\textbackslash ref} or
\texttt{\textbackslash eqref}. This chapter is
Chapter~\ref{ch:intro}, and the maths chapter is
Chapter~\ref{ch:maths}. Compile twice so the numbers settle.

The template also loads \emph{cleveref}, which adds the type word
for you: \cref{enum:start} needs no typed ``item'', and
\cref{ch:maths} comes out as a full ``chapter'' reference. Use
\texttt{\textbackslash cref} mid-sentence and \texttt{\textbackslash Cref} to open one.

\section{Citations and the index}
Cite anything in \texttt{references/reference.bib} by its key: the
classic advice is Knuth~\cite{knuth1984}, and modern practice is well
summarised elsewhere~\cite{lamport1994, mittal2005}. The full list is
printed automatically in the back matter, numbered in order of first
citation.

For the index, wrap any word in \texttt{\textbackslash index}: this book talks about
\index{typesetting}typesetting, \index{page breaks}page breaks and
the \index{index}index itself. Sub-entries use an exclamation mark:
\index{index!sub-entries}like so.

\section{A little longer prose}\label{sec:prose}
Good book typography is invisible: the reader should feel the page,
not notice it.\footnote{The classic reference here is
\textcite{bringhurst2004}, whose rules on measure, leading and rag
still hold on screen as on paper.} Three settings do most of that
work in this template --- measure (the line length set by the paper
format), leading (the line spacing in TYPOGRAPHY) and the colour of
the text block (how evenly the words grey the page).

The measure deserves respect. A line of running text wants between
sixty and seventy-five characters; shorter tires the eye with turns,
longer loses the return journey to the left margin. The margins in
each \texttt{\textbackslash GeomOpts} block are tuned with
that in mind, which is why they differ per paper size.

\section{Description lists}
A third list kind, for term--definition pairs:

\begin{description}[itemsep=2pt]
  \item[Measure] the width of the text block, usually given in ems or
    characters per line.
  \item[Leading] the baseline-to-baseline distance; the template's
    default is the font's own.
  \item[Folio] the printed page number, set on the outer edge here.
\end{description}

\noindent
The sources behind these terms, and everything else in this chapter,
are collected in the bibliography: see \parencite{knuth1984,
lamport1994} for the typesetting machinery and \textcite{texgyre}
for the fonts themselves.\index{leading}\index{measure|see{line length}}
